% docs/manual/user/chapters/01-overview.tex
% 第 1 章：概述与安装

\chapter{概述与安装}

\openbench{} 是一个面向陆面模型研究者的开源评估框架。本章介绍 \openbench{} 能做什么、为什么要用它、系统要求，以及在不同环境下的安装方法。

\section{\openbench{} 是什么}

\openbench{} 是一个全自动、跨平台的陆面模型基准测试框架。它把"用观测/再分析数据评估一个或多个陆面模型"这件原本繁琐、易错、难复现的事情，整合成一个 YAML 配置 + 一条命令的工作流。

具体来说，\openbench{} 解决三类问题：

\begin{itemize}
  \item \textbf{数据准备}：把分散在多个目录、多种分辨率、多种文件命名约定的 reference 与模型输出统一组织、对齐时间窗口、重网格化到目标分辨率；
  \item \textbf{评估计算}：用 25+ 种 metric 与归一化 score 评估每对 simulation × reference，支持 grid 与 station 两种数据形态；
  \item \textbf{结果呈现}：自动生成可视化、统计分析、HTML/PDF 报告，多模型对比按一致 layout 输出。
\end{itemize}

与同类工具（ILAMB、PLUMBER2 之类）相比，\openbench{} 的差异化在于：单一配置文件、内置的 reference 数据集注册表、对中文用户友好的 CLI/GUI、支持 SSH 远程执行、以及与 CoLM 系列模型的深度整合。

\begin{noteBox}
\openbench{} 不是数据下载器，也不是模型运行框架。它假设你已经有了 reference 数据集（或会通过 \cli{openbench data download} 让它帮你下载已注册的数据集）和一组陆面模型输出（NetCDF 文件）。
\end{noteBox}

\section{核心特性}

\begin{itemize}
  \item \textbf{统一配置}：单个 \file{openbench.yaml} 替代 v2.0 中需要的 4-6 个配置文件；
  \item \textbf{Reference 注册表}：内置 ~67 个常用数据集（GLEAM、ERA5LAND、FLUXNET、GRDC 等），按变量分类即查即用；
  \item \textbf{Model profile}：内置 21 种主流陆面模型（CoLM2024、CLM5、ERA5-Land 等）的变量映射模板，开箱即用；
  \item \textbf{完整 CLI}：\cli{openbench run / check / init / data / model / migrate / gui} 七大主命令；
  \item \textbf{可选 GUI}：14 步向导式配置生成器，适合不熟悉 YAML 的用户；
  \item \textbf{远程执行}：通过 SSH 在 HPC 集群上提交评估任务，本地查看进度；
  \item \textbf{迁移工具}：\cli{openbench migrate} 把 v2.0 JSON 或 v1.x Fortran NML 配置一键转换；
  \item \textbf{增量缓存}：基于 zarr 的中间结果缓存，重跑只重算变化部分。
\end{itemize}

\section{系统要求}

\begin{itemize}
  \item \textbf{Python}：3.10、3.11 或 3.12（不支持 3.9 及更早版本）
  \item \textbf{操作系统}：macOS、Linux（HPC 友好）；Windows 可用但未充分测试
  \item \textbf{内存}：最低 8 GB；典型全球 monthly × 5 年评估推荐 16 GB；千站点 hourly 评估推荐 32 GB+
  \item \textbf{磁盘}：reference 数据集本身可达数十到数百 GB；中间缓存通常占数据集体积的 0.5–2 倍
\end{itemize}

核心依赖（\openbench{} 自动管理，无需手动安装）：

\begin{itemize}
  \item 数据处理：\modname{xarray}、\modname{numpy}、\modname{scipy}、\modname{pandas}、\modname{netCDF4}、\modname{dask}、\modname{flox}
  \item 可视化：\modname{matplotlib}、\modname{cartopy}
  \item 配置与 CLI：\modname{PyYAML}、\modname{click}、\modname{jinja2}、\modname{platformdirs}
\end{itemize}

\section{安装}

\openbench{} 支持三种安装方式：\cli{pip}（最通用）、\cli{uv}（最快）、\cli{conda}（HPC 与离线友好，待发布到 conda-forge）。

\subsection{pip 安装}

最小安装（仅 CLI、HPC 友好）：

\begin{minted}{bash}
pip install openbench
\end{minted}

按需添加可选功能：

\begin{minted}{bash}
# 启用 GUI 向导
pip install "openbench[gui]"

# 启用 SSH 远程执行
pip install "openbench[remote]"

# 启用 PDF 报告生成
pip install "openbench[report]"

# 一次装齐所有可选功能
pip install "openbench[all]"
\end{minted}

\subsection{uv 安装}

\modname{uv} 是 Astral 开发的现代 Python 包管理器，比 pip 快 10-100 倍。如果你已经用 uv 管理项目：

\begin{minted}{bash}
uv pip install "openbench[all]"
\end{minted}

\subsection{conda 安装}

\warninline{conda-forge 通道暂未发布，本节为预留。}

发布后可用：

\begin{minted}{bash}
conda install -c conda-forge openbench
\end{minted}

\subsection{源码安装（开发者）}

如果你要修改代码或撰写贡献，从 GitHub 克隆并以"可编辑模式"安装：

\begin{minted}{bash}
git clone https://github.com/CoLM-SYSU/OpenBench.git
cd OpenBench
pip install -e ".[all,dev]"
\end{minted}

\cli{[dev]} 额外引入 \modname{pytest}、\modname{pytest-cov}、\modname{ruff} 用于测试与 lint。

\section{验证安装}

安装完成后，运行以下命令确认环境就绪：

\begin{minted}{bash}
openbench version
\end{minted}

预期输出：

\begin{minted}{text}
openbench 3.0.0a1
\end{minted}

查看可用命令：

\begin{minted}{bash}
openbench --help
\end{minted}

预期看到 7 个主命令（\cli{run}、\cli{check}、\cli{init}、\cli{data}、\cli{model}、\cli{migrate}、\cli{gui}）外加 \cli{version}。

\begin{tipBox}
\openbench{} CLI 使用 \textbf{lazy loading} ——首次执行某子命令时才加载对应模块。因此 \cli{openbench --help} 启动很快（< 0.5 秒），但首次 \cli{openbench run} 会有几百毫秒的额外加载时间。这是预期行为，不是 bug。
\end{tipBox}

\section{常见安装错误}

\subsection{cartopy 编译失败}

\modname{cartopy} 依赖系统级库（GEOS、PROJ）。若 \cli{pip install} 报编译错误，先安装系统依赖：

\begin{minted}{bash}
# macOS
brew install geos proj

# Ubuntu/Debian
sudo apt install libgeos-dev libproj-dev proj-data proj-bin
\end{minted}

或者用 conda 安装 cartopy 后再用 pip 装 \openbench{}：

\begin{minted}{bash}
conda install -c conda-forge cartopy
pip install openbench
\end{minted}

\subsection{PySide6 在 HPC 上装不上}

\cli{[gui]} 引入 \modname{PySide6}（Qt 6 绑定），在无图形系统的 HPC 节点上往往装不上。在 HPC 上不要装 \cli{[gui]}：

\begin{minted}{bash}
# HPC 节点典型安装
pip install "openbench[remote]"  # 不带 gui
\end{minted}

GUI 配置生成本来就该在本地做完再上传配置文件。

\subsection{在受限网络上离线安装}

把 \openbench{} 及全部依赖打包成 wheelhouse：

\begin{minted}{bash}
# 在有网络的机器上
pip download openbench[all] -d ./wheelhouse

# 拷贝 wheelhouse 到目标机器后
pip install --no-index --find-links=./wheelhouse openbench[all]
\end{minted}

更详细的 HPC 部署见运维卷第 2 章。

\section{下一步}

环境就绪后，下一章会用 10 分钟带你跑通第一个完整评估，包括：

\begin{itemize}
  \item 用 \cli{openbench init} 交互式生成最小配置
  \item 用 \cli{openbench check} 验证配置与数据
  \item 用 \cli{openbench run} 启动评估
  \item 看输出目录里有什么
\end{itemize}
